% ---- pg-567 （书页 555） ----

\begin{mvfigure}
\begin{minipage}[t]{0.47\linewidth}\centering
\begin{tikzpicture}[x=0.68cm,y=28cm,font=\zihao{6}]
\draw[->] (0,-0.006) -- (6.6,0) node[below right,inner sep=1pt] {$q$};
\draw[->] (0,-0.006) -- (0,0.152);
\node[rotate=90,inner sep=1pt] at (-0.72,0.062) {应力};
\foreach \x in {2,4,6} {\draw (\x,0) -- (\x,-0.005);}
\node[below,inner sep=1.5pt] at (2,-0.005) {2};
\node[below,inner sep=1.5pt] at (4,-0.005) {4};
\node[below,inner sep=1.5pt] at (6,-0.005) {6};
\foreach \y in {0,0.02,0.04,0.06,0.08,0.10,0.12,0.14} {\draw (0,\y) -- (-0.22,\y);}
\node[left,inner sep=1.5pt] at (-0.22,0) {0};
\node[left,inner sep=1.5pt] at (-0.22,0.02) {0.02};
\node[left,inner sep=1.5pt] at (-0.22,0.04) {0.04};
\node[left,inner sep=1.5pt] at (-0.22,0.06) {0.06};
\node[left,inner sep=1.5pt] at (-0.22,0.08) {0.08};
\node[left,inner sep=1.5pt] at (-0.22,0.10) {0.10};
\node[left,inner sep=1.5pt] at (-0.22,0.12) {0.12};
\node[left,inner sep=1.5pt] at (-0.22,0.14) {0.14};
\draw (1,0.125) -- (2,0.007) -- (3,0.013) -- (4,0.008);
\fill (1,0.125) circle (0.9pt);
\fill (2,0.007) circle (0.9pt);
\fill (3,0.013) circle (0.9pt);
\fill (4,0.008) circle (0.9pt);
\node[above,inner sep=1.5pt] at (1,0.125) {1};
\node[below left,inner sep=1.5pt] at (2,0.007) {2};
\node[above,inner sep=1.5pt] at (3,0.013) {3};
\node[right,inner sep=1.5pt] at (4,0.008) {4};
\end{tikzpicture}
\phantomsection\label{fig:12.16}\mvcap{图 12.16\quad 城市间航线距离的应力函数}
\end{minipage}\hfill
\begin{minipage}[t]{0.47\linewidth}\centering
\begin{tikzpicture}[x=0.52cm,y=9.5cm,font=\zihao{6}]
\draw[->] (0,-0.055) -- (9.0,0) node[below right,inner sep=1pt] {$q$};
\draw[->] (0,-0.055) -- (0,0.415);
\node[rotate=90,inner sep=1pt] at (-0.75,0.15) {应力};
\foreach \x in {0,2,4,6,8} {\draw (\x,-0.055) -- (\x,-0.075);}
\node[below,inner sep=1.5pt] at (0,-0.075) {0};
\node[below,inner sep=1.5pt] at (2,-0.075) {2};
\node[below,inner sep=1.5pt] at (4,-0.075) {4};
\node[below,inner sep=1.5pt] at (6,-0.075) {6};
\node[below,inner sep=1.5pt] at (8,-0.075) {8};
\foreach \y in {-0.05,0.00,0.05,0.10,0.15,0.20,0.25,0.30,0.35,0.40} {\draw (0,\y) -- (-0.4,\y);}
\node[left,inner sep=1.5pt] at (-0.4,-0.05) {$-0.05$};
\node[left,inner sep=1.5pt] at (-0.4,0.00) {$0.00$};
\node[left,inner sep=1.5pt] at (-0.4,0.05) {$0.05$};
\node[left,inner sep=1.5pt] at (-0.4,0.10) {$0.10$};
\node[left,inner sep=1.5pt] at (-0.4,0.15) {$0.15$};
\node[left,inner sep=1.5pt] at (-0.4,0.20) {$0.20$};
\node[left,inner sep=1.5pt] at (-0.4,0.25) {$0.25$};
\node[left,inner sep=1.5pt] at (-0.4,0.30) {$0.30$};
\node[left,inner sep=1.5pt] at (-0.4,0.35) {$0.35$};
\node[left,inner sep=1.5pt] at (-0.4,0.40) {$0.40$};
\draw (1,0.375) -- (2,0.165) -- (3,0.072);
\draw[dashed] (3,0.072) -- (4,0.025) -- (5,0.008) -- (6,0.010);
\fill (1,0.375) circle (0.9pt);
\fill (2,0.165) circle (0.9pt);
\fill (3,0.072) circle (0.9pt);
\fill (4,0.025) circle (0.9pt);
\fill (5,0.008) circle (0.9pt);
\fill (6,0.010) circle (0.9pt);
\node[above,inner sep=1.5pt] at (1,0.375) {1};
\node[left,inner sep=1.5pt] at (2,0.165) {2};
\node[below,inner sep=1.5pt] at (3,0.072) {3};
\node[below,inner sep=1.5pt] at (4,0.025) {4};
\node[below,inner sep=1.5pt] at (5,0.008) {5};
\node[right,inner sep=1.5pt] at (6,0.010) {6};
\end{tikzpicture}
\phantomsection\label{fig:12.17}\mvcap{图 12.17\quad 公用事业公司间距离的应力函数}
\end{minipage}
\end{mvfigure}

\hyperref[fig:12.17]{图 12.17} 中的应力函数没有明显的肘形. 此图显示“好”应力值(小于或等于 $5\%$)的点大致在 $q=4$ 附近. 对公用事业公司可以得出一个好的 4 维表示, 但很难作图. 我们在\hyperref[fig:12.18]{图 12.18} 中展示 $q=2$ 维的公用事业公司结构图. 坐标轴沿最终散点的样本主成分方向.

\begin{mvfigure}
\mvinc{pg-567}{0.240,0.504,0.779,0.741}
\phantomsection\label{fig:12.18}\mvcap{图 12.18\quad 由多维标度变换产生的公用事业公司的几何表示}
\end{mvfigure}

虽然二维情形下的应力值相当高(应力$(2)\times 100\%=19\%$), \hyperref[fig:12.18]{图 12.18} 中诸公司间的距离与本章早些时候所提供的聚类分析结果差别不是特别大. 例如, 中西部的公司——联邦爱迪生、威斯康星电力(WEPCO)、麦迪逊天然气与电力(MG\&E), 以及北方电力(NSP)——彼此很接近(相似). 得克萨斯公用事业及俄克拉何马天然气与电力(Ok. G.\&E)也非常接近(相似). 其他公用事业公司倾向于按地理位置或相似的环境分组.
